\section{Technology Stack}
\begin{enumerate}[label=\Alph*., leftmargin=1.5em]
    \item Frontend: HTML5, CSS3, JavaScript
    \item Backend: Node.js, Express.js
    \item Database: MySQL
    \item Security: JWT, bcryptjs
    \item Email: Nodemailer
    \item Visualization: Chart.js
\end{enumerate}

\section{Development Environment}
The project was developed using Visual Studio Code, the Node.js runtime, npm for dependency management, and MySQL database tooling. Version control was maintained using Git during feature integration.

\section{Backend Implementation}
The backend is implemented using Express.js with modular routes and middleware. Core implementation points include:
\begin{enumerate}[label=\Alph*., leftmargin=1.5em]
    \item secure login and registration endpoints,
    \item JWT token generation and verification through auth middleware,
    \item performance APIs for teacher ranking and trend retrieval,
    \item admin-only endpoints for individual and bulk email operations.
\end{enumerate}
Database access is handled through MySQL query interfaces, and password security is enforced using bcrypt hashing.

\section{Frontend Implementation}
The frontend is organized into role-aware dashboards and reusable UI sections. Implementation includes:
\begin{enumerate}[label=\Alph*., leftmargin=1.5em]
    \item authentication pages for login and registration,
    \item dashboard tabs for home, professor list, doubts, and analytics,
    \item rating and feedback forms with validation,
    \item performance prediction modal with multi-step input,
    \item chart rendering using Chart.js for trend visualization.
\end{enumerate}

\section{Prediction Logic}
The performance module uses a hybrid strategy:
\begin{enumerate}[label=\Alph*., leftmargin=1.5em]
    \item weighted formula based on study behavior and previous performance,
    \item regression blending using historical entries for better estimation,
    \item final output in the format \textbf{Expected Score: X / 10}.
\end{enumerate}

\section{Error Handling and Validation}
Input validation is implemented at both frontend and backend levels to prevent invalid records. API responses use structured status and message handling to provide clearer user feedback and easier debugging.

\section{Limitations}
The current implementation is optimized for institute-scale deployment; however, it can be further enhanced through advanced model calibration, centralized logging, and deeper historical analytics.

\section*{Chapter Summary}
This chapter described the concrete implementation of the EduRate platform across frontend, backend, and data layers. It also summarized technical choices, validation strategy, prediction logic, and current system limitations for future enhancement planning.
